\subsection{External hilbertian sum of hilbertian spaces}

PROPOSITION 1. --- Let \((\mathrm{E}_i)_{i\in \mathbf{I}}\) be a family of hilbertian spaces, \(\mathbf{P}\) the product vector space \(\prod_{i\in \mathbf{I}}\mathrm{E}_i\) and \(\mathbf{E}\) the subset of \(\mathbf{P}\) consisting of all families \(\mathbf{x} = (x_i)_{i\in \mathbf{I}}\) such that \(\sum_{i\in \mathbf{I}}\| x_i\|^2\) is finite.

\begin{enumerate}
\def\labelenumi{\alph{enumi})}
\item
  E is a vector subspace of \(\mathbf{P}\) .
\item
  For every \(\mathbf{x} = (x_i)_{i\in \mathbf{I}}\) and \(\mathbf{y} = (y_i)_{i\in \mathbf{I}}\) in E, the family \((\langle x_i|y_i\rangle)_{i\in \mathbf{I}}\) is summable. If we put \(\langle \mathbf{x}|\mathbf{y}\rangle = \sum_{i\in \mathbf{I}}\langle x_i|y_i\rangle\) , we define a positive separating hermitian form on E.
\item
  For the scalar product so defined, \(\mathbf{E}\) is a hilbertian space; the direct sum \(\mathbf{S}\) of the \(\mathbf{E}_i\) is dense in \(\mathbf{E}\) .
\end{enumerate}

For \(\mathbf{x} = (x_{i})_{i\in \mathbf{I}}\) and \(\mathbf{y} = (y_i)_{i\in \mathbf{I}}\) in E, we have

\[
\left\| x _ {i} + y _ {i} \right\| ^ {2} \leqslant 2 \left(\left\| x _ {i} \right\| ^ {2} + \left\| y _ {i} \right\| ^ {2}\right),
\]

hence \(\mathbf{x} + \mathbf{y} = (x_i + y_i)_{i\in I}\) belongs to E. This proves \(a)\) .

By the Cauchy-Schwarz inequality, we have

\[
\left| \left\langle x _ {i} \mid y _ {i} \right\rangle \right| \leqslant \| x _ {i} \| \cdot \| y _ {i} \| \leqslant \frac {1}{2} \left(\left\| x _ {i} \right\| ^ {2} + \left\| y _ {i} \right\| ^ {2}\right)
\]

hence \(\sum_{i\in \mathbf{I}}|\langle x_i|y_i\rangle | < +\infty\) . If \(x\neq 0\) , we have \(\langle \mathbf{x}|\mathbf{x}\rangle = \sum_{i\in \mathbf{I}}\| x_i\|^2 >0\) , hence assertion \(b)\) follows.

We recall that S is the subspace of P consisting of all families \(\mathbf{x} = (x_{i})_{i \in \mathbf{I}}\) such that the set of all \(i \in I\) for which \(x_{i} \neq 0\) is finite. It follows immediately that S is dense in E; hence it remains to prove that E is complete for the topology \(T_{1}\) obtained by the norm \(\|x\| = \langle x|x\rangle^{1/2}\) . Let \(T_{2}\) be the topology induced on E by the product topology on \(\prod_{i \in I} E_{i}\) . For every r \textgreater{} 0, let \(B_{r}\) be the set of all \(x \in E\) such that \(\|x\| \leqslant r\) . This relation implies that we have \(\sum_{i \in J} \|x_{i}\|^{2} \leqslant r^{2}\) for every finite subset J of I, and so \(B_{r}\) is a closed subset of \(\prod_{i \in I} E_{i}\) , hence also complete. The fact that E is complete for \(T_{1}\) now follows from GT, III, §3, No.~5, cor. 2 to prop. 10.

DEFINITION 1. --- Let \((\mathrm{E}_{i})_{i\in\mathrm{I}}\) be a family of hilbertian spaces. The hilbertian space E defined in prop. 1 is called the external hilbertian sum of the family \((\mathrm{E}_{i})_{i\in\mathrm{I}}\) and written as \(\bigoplus_{i\in I} E_{i}\) or \(\bigoplus_{i\in I} E_{i}^{1}\) .

\(^{1}\) Care must be taken not to confuse this notation with that of the « algebraic » direct sum of the spaces \(E_{i}\) (A, II, § 1, No.~6).

Let \(f_{i}\) be the mapping from \(E_{i}\) into \(E\) which transforms \(z \in E_{i}\) into an element \((x_{k}) \in E\) such that \(x_{k} = 0\) for all \(k \neq i\) and \(x_{i} = z\) ; it is clear that \(f_{i}\) is an isomorphism from the hilbertian space \(E_{i}\) onto a closed vector subspace of \(E\) . We say that \(f_{i}\) is the canonical mapping from \(E_{i}\) into \(E\) and we shall generally identify \(E_{i}\) with its image in \(E\) by this isomorphism. With this convention, \(E_{i}\) and \(E_{k}\) are orthogonal in \(E\) for \(i \neq k\) , and \(E\) is the closed vector subspace generated by the union of the subspaces \(E_{i}\) .

When I is finite, E is the direct sum of the \(E_{i}\) ; since the canonical projector from E onto \(E_{i}\) is continuous for all \(i \in I\) , E is also the topological direct sum of the \(E_{i}\) (GT, III, § 6, No.~2, prop. 2). If \(I = \left(1, n\right)\) , we also write \(E_{1} \oplus E_{2} \oplus \ldots \oplus E_{n}\) instead of \(\bigoplus_{i=1}^{n} E_{i}\) .

Example. --- Let E be a hilbertian space and I a set of indices. Let \(\ell_{\mathrm{E}}^{2}(\mathrm{I})\) denote the external hilbertian sum of the family \((\mathrm{E}_{i})_{i\in\mathrm{I}}\) where \(E_{i}=E\) for all \(i\in I\) . In other words, \(\ell_{\mathrm{E}}^{2}(\mathrm{I})\) is the space of all families \(\mathbf{x}=(x_{i})_{i\in\mathrm{I}}\) of elements of E such that \(\sum_{i\in I}\|x_{i}\|^{2}<+\infty\) , endowed with the scalar product \(\langle x|y\rangle=\sum_{i\in I}\langle x_{i}|y_{i}\rangle\) (space of square summable families of elements of E indexed by I). We put \(\ell^{2}(\mathrm{I})=\ell_{\mathrm{K}}^{2}(\mathrm{I})\) .

